\documentclass[12pt]{article}
\usepackage[utf8]{inputenc}
\usepackage[T1]{fontenc}
\usepackage{hyperref}
\title{Chambri Language Resources for K-12 Curriculum Development}
\author{Compiled by Claude Sonnet 5.5 (claude-sonnet-5-5) for Andr\'as Kornai}
\date{}
\begin{document}
\maketitle

\section{Introduction}
Chambri (ISO 639-3 can; Glottocode cham1313; older spellings Tchambuli, Tshamberi) is a Lower Sepik (Ramu--Lower Sepik) language of the Chambri Lakes area, East Sepik Province, Papua New Guinea. Wikipedia gives about 800 speakers (2006), well below the 10,000 in the request; the Endangered Languages Project lists it as threatened with low certainty. Search found no published grammar or dictionary.

\section{Primary Educational Resources}
No primers, textbooks or school programmes for Chambri were found in search results.

\section{Linguistic and Reference Materials}
\begin{itemize}
\item Wikipedia, Chambri language (cites Foley's 2018 overview of Sepik-Ramu languages): \url{https://en.wikipedia.org/wiki/Chambri_language}
\item Chambri numeral system (Lean, 1991 data), MPI-EVA: \url{https://lingweb.eva.mpg.de/channumerals/Chambri.htm}
\item Foley (2005) comparative Lower Sepik lexical data including Chambri: \url{https://transnewguinea.org/source/foley2005?page=5}
\end{itemize}

\section{Religious and Community Resources}
None found in search results.

\section{Digital and Audio Resources}
\begin{itemize}
\item Endangered Languages Project entry: \url{https://endangeredlanguages.com/node/98275}
\end{itemize}
Search summaries mention Foley and Laycock recordings and notebooks in the Paradisec archive, but no Paradisec URL appeared in the results, so none is cited here.

\section{Curriculum Development Context}
With few documented resources and a small speaker base, curriculum work would need new documentation with the community, possibly starting from archived recordings. Elementary-school vernacular education in Papua New Guinea would be the natural institutional context.
\end{document}